\NativeTitle{85}{总复习（二）}{（图形与几何）}{90}{5}{一 & 二 & 三 & 四 & 五}{ &  &  &  & }
\NativeSlot{272}{40}{%
\ExamGroup{一、填空。{\normalfont\kaishu\fontsize{10.8bp}{14bp}\selectfont （每空1分，共23分）}}

% p085_s01

\ExamQuestion{1.}{5}%
一条直线的平行线有\mbox{（\hspace{1.8em}）}条，过直线外一点可以画\mbox{（\hspace{1.8em}）}条直线与这条直线垂直。\par
}
\NativeSlot{346}{36}{%
% p085_s02

\ExamQuestion{2.}{5}%
长方形的邻边相互\mbox{（\hspace{1.8em}）}，对边相互\mbox{（\hspace{1.8em}）}。\par
}
\NativeSlot{370}{36}{%
% p085_s03

\ExamQuestion{3.}{5}%
以平行四边形的一条边为底，能作出\mbox{（\hspace{1.8em}）}条高，这些高的长度\mbox{（\hspace{1.8em}）}。\par
}
\NativeSlot{395}{57}{%
\par\noindent%
\ExamTile{.74}{%
% p085_s04

\ExamQuestion{4.}{5}%
右图中，有\mbox{（\hspace{1.8em}）}个三角形，有\mbox{（\hspace{1.8em}）}个平行四边形，有\mbox{（\hspace{1.8em}）}个梯形。\par
}\hfill%
\ExamTile{.235}{%
\BookFigure{figures/p085_s04-01.tex}
}\par
% p085_s05

\ExamQuestion{5.}{5}%
已知一个三角形的两条边长分别是\(5\,\mathrm{cm}\)和\(9\,\mathrm{cm}\)，第三条边最长是\mbox{（\hspace{1.8em}）}\(\mathrm{cm}\)，最短是\mbox{（\hspace{1.8em}）}\(\mathrm{cm}\)。（取整厘米数）\par
}
\NativeSlot{495}{104}{%
\par\noindent%
\ExamTile{.74}{%
% p085_s06

\ExamQuestion{6.}{5}%
如图，在直角三角形\(ABC\)中（单位：厘米），\allowbreak{}\(BC\)边上的高是\mbox{（\hspace{1.8em}）}厘米，\allowbreak{}\(AB\)边上的高是\mbox{（\hspace{1.8em}）}厘米。\par
}\hfill%
\ExamTile{.235}{%
\BookFigure{figures/p085_s06-01.tex}
}\par
% p085_s07

\ExamQuestion{7.}{5}%
下面哪些是平移现象，哪些是旋转现象？\par\noindent （1）张叔叔在公路上开车，方向盘的运动是\mbox{（\hspace{1.8em}）}现象。\par

% p085_s08

\ExamQuestion{（2）}{5}%
升国旗时，国旗的上升运动是\mbox{（\hspace{1.8em}）}现象。\par

% p085_s09

\ExamQuestion{（3）}{5}%
妈妈用拖布拖地，是\mbox{（\hspace{1.8em}）}现象。\par

% p085_s10

\ExamQuestion{（4）}{5}%
小明骑自行车，车轮的运动是\mbox{（\hspace{1.8em}）}现象。\par
}
\NativeSlot{671}{36}{%
% p085_s11

\ExamQuestion{8.}{5}%
正方形有\mbox{（\hspace{1.8em}）}条对称轴，等边三角形有\mbox{（\hspace{1.8em}）}条对称轴，等腰梯形有\mbox{（\hspace{1.8em}）}条对称轴。\par
}
\NativeSlot{721}{37}{%
\par\noindent%
\ExamTile{.13}{%
\noindent\makebox[\linewidth][c]{\NativeImage{figures/p085_s12-01.tex}}
}\hfill%
\ExamTile{.845}{%
% p085_s12

\ExamQuestion{9.}{5}%
钟面显示的是\mbox{（\hspace{1.8em}）}时，将分针顺时针旋转一个直角后是\mbox{（\hspace{1.8em}）}时\mbox{（\hspace{1.8em}）}分。\par
}\par
}
